\documentclass[11pt]{article}
\usepackage[T1]{fontenc}
\usepackage[utf8]{inputenc}
\usepackage{lmodern}
\usepackage[margin=1in]{geometry}
\usepackage{amsmath,amssymb,amsthm,mathtools}
\usepackage{booktabs,tabularx,array}
\usepackage{enumitem}
\usepackage{aliascnt}
\usepackage{microtype}
\usepackage{xcolor}
\usepackage{hyperref}
\usepackage[nameinlink,noabbrev]{cleveref}

\definecolor{linkblue}{RGB}{28,65,113}
\hypersetup{colorlinks=true,linkcolor=linkblue,citecolor=linkblue,urlcolor=linkblue,
 pdftitle={Minimax Regret under Per-Round Separation-Oracle Budgets},
 pdfauthor={Colombo Team},pdfkeywords={online convex optimization, separation oracle, minimax regret, asphericity}}
\setlist[itemize]{leftmargin=1.5em,itemsep=3pt,topsep=4pt}
\setlist[enumerate]{leftmargin=1.8em,itemsep=3pt,topsep=4pt}
\setlength{\emergencystretch}{2em}
\setlength{\belowcaptionskip}{5pt}
\newcolumntype{Y}{>{\raggedright\arraybackslash}X}
\newcolumntype{P}[1]{>{\raggedright\arraybackslash}p{#1}}
\newcommand{\R}{\mathbb R}
\newcommand{\E}{\mathbb E}
\newcommand{\Prob}{\mathbb P}
\newcommand{\K}{\mathcal K}
\newcommand{\Reg}{\operatorname{Reg}}
\newcommand{\vol}{\operatorname{vol}}
\newcommand{\Cov}{\operatorname{Cov}}
\newcommand{\Sep}{\mathsf{Sep}}
\newcommand{\Back}{\mathsf{Back}}
\newcommand{\ip}[2]{\langle #1,#2\rangle}
\newcommand{\norm}[1]{\lVert #1\rVert}
\newcommand{\wtO}{\widetilde O}
\newcommand{\wtTheta}{\widetilde\Theta}
\newcommand{\eps}{\varepsilon}
\newcommand{\ind}{\mathbf 1}
\DeclareMathOperator{\proj}{proj}

\newtheorem{theorem}{Theorem}[section]
\newaliascnt{lemma}{theorem}
\newtheorem{lemma}[lemma]{Lemma}
\aliascntresetthe{lemma}
\newaliascnt{corollary}{theorem}
\newtheorem{corollary}[corollary]{Corollary}
\aliascntresetthe{corollary}
\theoremstyle{definition}
\newaliascnt{definition}{theorem}
\newtheorem{definition}[definition]{Definition}
\aliascntresetthe{definition}
\newcounter{algorithm}
\crefname{algorithm}{Algorithm}{Algorithms}
\Crefname{algorithm}{Algorithm}{Algorithms}

\title{Minimax Regret under Per-Round\\Separation-Oracle Budgets}
\author{Colombo Team\\\normalsize\href{mailto:colombo.2026@outlook.com}{\texttt{colombo.2026@outlook.com}}}
\date{}

\begin{document}
\maketitle

\begin{abstract}
We study adversarial online convex optimization over an unknown convex body
$rB_2^d\subseteq\K\subseteq RB_2^d$ accessed through an exact separation oracle.
The learner must play feasible points and obey an oracle budget on every round.
Writing $\kappa=R/r$ and $M=\min\{T,d\}$, we show that the minimax expected regret
under a sufficiently large, strictly enforced polylogarithmic budget is
\[
 \wtTheta\!\left(GR\min\{T,\sqrt T+\min\{M,\kappa\sqrt M\}\}\right),
\]
where $G$ bounds the subgradient norm and logarithms in $d,T,\kappa$ are
suppressed. The upper bound combines Gaussian localization with a deterministic
separation-based learner on rounds when localization fails. Strong
log-concavity controls online stability without a dimension factor, while each
rejected approximate centroid supplies a constant decrease in a partition
potential. Both components share a single per-round query budget. The lower
bound hides independent signs in the facets of a ball. Every adaptive oracle
transcript leaves a uniform posterior on a coordinate subcube, and one oracle
answer identifies at most two new signs in conditional expectation. Exact
feasibility then forces regret until the hidden facets have been learned.
We also give explicit bounds for a general query cap $Q$. These bounds match up
to logarithms for polylogarithmic $Q$, but leave an intermediate-asphericity gap
for larger budgets. The results concern unknown-body oracle complexity; the
centroid implementation does not provide a competitive arithmetic-time bound.
\end{abstract}

\noindent\textbf{Keywords.} Online convex optimization; separation oracles;
minimax regret; asphericity; oracle complexity; log-concave localization.

\medskip
\begin{samepage}
\noindent\textbf{AI Usage.}
Nearly the entire research pipeline for this paper was carried out by
Colombo, powered by GPT-5.6 Sol. 
\par
\end{samepage}

\section{Introduction}

Online convex optimization asks a learner to select feasible decisions before
observing their losses. With access to Euclidean projections, online gradient
descent achieves regret $O(GR\sqrt T)$ over a convex set contained in a ball of
radius $R$~\cite{zinkevich2003}. For a set specified implicitly, however, a
projection may be substantially more expensive than a feasibility query. A
separation oracle provides an alternative: it either certifies that a proposed
point is feasible or returns a hyperplane separating that point from the set.
The resulting question is how much regret must be paid when only a small
number of such queries is available before each decision.

Two features distinguish this question from the usual analysis of online
gradient descent. First, the learner must discover useful parts of the
feasible set while learning from the losses. Second, a budget imposed on each
round prevents the learner from spending an arbitrarily large number of
queries before its first decision. A linear total query budget does not impose
this restriction. The geometry of the set also matters: if a known inner ball
has radius $r$, then the ratio $\kappa=R/r$ measures the geometric uncertainty
allowed by the prescribed ball sandwich.

We characterize the interaction between horizon, dimension, and this
asphericity bound under a strict polylogarithmic query budget. With
$M=\min\{T,d\}$, the answer is
\begin{equation}\label{eq:headline}
 \wtTheta\!\left(GR\min\{T,\sqrt T+\min\{M,\kappa\sqrt M\}\}\right).
\end{equation}
The term $\sqrt T$ is the unavoidable cost of online learning even when the
body is known. The remaining term describes the cost of learning an unknown
body. It grows with $\kappa$ until it saturates at $M$. This geometric cost
need not be lower order: when the horizon is short relative to the dimension,
it can make regret linear in $T$. Conversely, if $M\le\sqrt T$, the standard
$\sqrt T$ rate is attainable for every asphericity, up to logarithms.

\paragraph{Contributions.}
The analysis has three components.
\begin{enumerate}
 \item \textbf{An upper bound that separates stability from geometry.}
 A truncated Gaussian over an explicit outer approximation of $\K$ gives
 dimension-independent online stability. A rejected approximate centroid
 reduces the log partition function by a constant amount. Charging these
 decreases against the regret of a deterministic fallback learner yields
 an upper bound of order
 $GR[\sqrt T+\kappa\sqrt{d/q}]$, up to logarithms, where $q$ is the number
 of localization queries reserved on each round. An origin fallback gives
 $\wtO(GR[\sqrt T+d/Q])$ with any cap $Q\ge1$.
 \item \textbf{A lower bound that tracks both queries and asphericity.}
 A family of hidden-facet bodies gives an explicit comparator advantage for
 every $\kappa>1$. An exact posterior invariant bounds the information
 revealed by adaptive separation queries. The argument applies to randomized
 learners required to remain feasible almost surely on every admissible body.
 \item \textbf{A minimax characterization with an explicit resource model.}
 Combining the bounds proves \eqref{eq:headline} for, in particular, the
 logarithmic cap $Q_\star=2[1+\lceil\log_2(\kappa T)\rceil]$.
 The same order holds for larger caps bounded by a fixed power of a logarithm.
 For general growing $Q$, we state the remaining gap explicitly.
\end{enumerate}

The deterministic fallback uses the established gauge-reduction idea of
Mhammedi~\cite{mhammedi2022membership,mhammedi2025sep}; we include a complete
proof with its precise query cost. Our contribution is the localization
analysis, its combination with that reduction, and the matching geometric
lower bound. The distinction matters because a membership-only implementation
of the same regret guarantee has a different query cost. A comparison of
the relevant models and bounds appears in \cref{sec:related}.

\section{Model and main results}\label{sec:results}

\subsection{Protocol and charged resources}

Let $d,T\ge1$ be integers and let $G,R,r>0$, with $r\le R$.
The learner knows these parameters and the common center, taken to be the
origin. The unknown compact convex body satisfies
\begin{equation}\label{eq:sandwich}
 rB_2^d\subseteq\K\subseteq RB_2^d,
 \qquad \kappa:=R/r.
\end{equation}
Here $B_2^d$ is the closed Euclidean unit ball. The supplied radii need not be
optimal for $\K$; thus $\kappa$ is a bound associated with the model class.
In particular, the known ball $RB_2^d$ belongs to this class for every
$\kappa\ge1$.

At round $t$, the learner first makes at most $Q$ oracle calls and then plays
$w_t\in\K$. It receives a convex loss $f_t$ and a subgradient
$g_t\in\partial f_t(w_t)$ satisfying $\norm{g_t}_2\le G$.
The loss may depend on the history and the current action. The algorithms
below use only $g_t$. Their regret is
\begin{equation}\label{eq:regret}
 \Reg_T:=\sup_{u\in\K}\sum_{t=1}^T[f_t(w_t)-f_t(u)]
 \le \sup_{u\in\K}\sum_{t=1}^T\ip{g_t}{w_t-u}.
\end{equation}
The upper bounds hold uniformly over comparators and realized subgradient
sequences on the event of accurate centroid computation.

\begin{definition}[Exact strong separation]\label{def:sep}
For a query $z\in\R^d$, the oracle $\Sep_\K$ returns either
\texttt{inside}, certifying $z\in\K$, or a unit vector $a$ such that
\[
 \sup_{x\in\K}\ip{a}{x}<\ip{a}{z}.
\]
In the latter case, $\{x:\ip{a}{x}\le\ip{a}{z}\}$ is a valid outer halfspace.
The oracle may select any valid response.
\end{definition}

The cap includes every call made before the current action, starting on
round~1. No uncharged preprocessing access to $\K$ is allowed. Calls cannot
be saved and spent on later rounds. Internal arithmetic and queries to
explicitly stored inequalities are not charged as unknown-body oracle calls.
We impose no arithmetic-time or bit-complexity bound. \Cref{app:sampling}
gives an explicit sampling implementation with finite expected internal work;
this work can be large.

Write $\mathfrak R_Q^{\rm sep}(d,T,\kappa)$ for the infimum, over learners
obeying this cap and feasible almost surely for every body and valid oracle,
of the worst-case expected regret. The supremum ranges over bodies satisfying
\eqref{eq:sandwich}, valid separation responses, and admissible loss
adversaries. Dependence on $G,R$ is suppressed. All constants below are
universal. Tilde notation suppresses factors polynomial in
$\log(edT\kappa)$.

\subsection{Upper bounds for a fixed cap}

Define the explicit quantities
\begin{equation}\label{eq:LD}
 L_T:=1+\lceil\log_2(\kappa T)\rceil,
 \qquad D_T:=d\log(e\kappa T).
\end{equation}
The integer $L_T$ reserves enough separation queries for one action of the
deterministic fallback learner in \cref{sec:backup}.

\begin{theorem}[Upper bounds]\label{thm:upper}
For every integer $Q\ge1$, an almost surely feasible randomized learner using
at most $Q$ exact separation calls on every round satisfies
\begin{equation}\label{eq:upper-origin}
 \mathfrak R_Q^{\rm sep}
 \le C GR\min\left\{T,\sqrt T+\frac{D_T}{Q}\right\}.
\end{equation}
If $Q\ge L_T+1$ and $q:=Q-L_T$, then also
\begin{equation}\label{eq:upper-hybrid}
 \mathfrak R_Q^{\rm sep}
 \le C GR\min\left\{
 T,\ \kappa\sqrt T,\ \sqrt T+\frac{D_T}{Q},\
 \sqrt T+\kappa\sqrt{\frac{D_T}{q}}
 \right\}.
\end{equation}
The bounds hold against adaptive loss adversaries.
\end{theorem}

The minimum is implemented by selecting a strategy using the known
parameters, rather than by querying all strategies in parallel. The
origin-fallback strategy uses the entire budget $Q$ for localization.
The hybrid reserves $L_T$ calls for its fallback and uses the remaining
$q=Q-L_T$ for localization. These are separate upper-bound constructions.

\subsection{Lower bounds for a fixed cap}

For $m\ge1$ and $\kappa\ge1$, set
\begin{equation}\label{eq:delta}
 \delta_{m,\kappa}:=
 \frac{\kappa^2-1}
 {2\kappa^2\left(
 \sqrt{(\kappa^2+1)/(2\kappa^2)}+\sqrt{1+m/\kappa^2}
 \right)}.
\end{equation}
This quantity is zero at $\kappa=1$ and is the normalized comparator
advantage of the hidden-facet construction.

\begin{theorem}[Lower bounds]\label{thm:lower}
For $d\ge2$, $T\ge1$, $\kappa\ge1$, and every integer $Q\ge1$,
\begin{equation}\label{eq:lower-exact}
 \mathfrak R_Q^{\rm sep}
 \ge cGR\max\left\{\sqrt T,
 \max_{1\le m\le d-1}
 \delta_{m,\kappa}
 \left\lfloor\min\left\{T,\frac{m}{4Q}\right\}\right\rfloor
 \right\}.
\end{equation}
For $\kappa\ge2$, with $N_Q:=\min\{T,(d-1)/Q\}$, this implies
\begin{equation}\label{eq:lower-simple}
 \mathfrak R_Q^{\rm sep}
 \ge cGR\min\left\{T,
 \sqrt T+\min\left\{N_Q,\kappa\sqrt{\frac{N_Q}{Q}}\right\}
 \right\}.
\end{equation}
The hidden-facet bound uses one fixed legal oracle for each body and the
same linear loss on every round. The $\sqrt T$ bound already holds on a
known ball.
\end{theorem}

The floor in \eqref{eq:lower-exact} retains the small-budget and
small-dimension endpoints literally. Its effect is absorbed by the
$\sqrt T$ term only when passing to \eqref{eq:lower-simple}.

\subsection{Minimax rate with a strict polylogarithmic budget}

To make the budget convention unambiguous, define
\begin{equation}\label{eq:pldef}
 Q_\star:=2L_T,
 \qquad
 \mathfrak R_{\rm pl}^{\rm sep}:=\mathfrak R_{Q_\star}^{\rm sep}.
\end{equation}
The notation does not refer to every cap bounded above by a polylogarithm;
in particular, it does not assert the hybrid result for $Q=1$.

\begin{corollary}[All-asphericity minimax rate]\label{cor:main}
For $d,T\ge1$, $\kappa\ge1$, and $M:=\min\{T,d\}$,
\begin{equation}\label{eq:main}
 \mathfrak R_{\rm pl}^{\rm sep}(d,T,\kappa)
 =\wtTheta\!\left(
 GR\min\{T,\sqrt T+\min\{M,\kappa\sqrt M\}\}
 \right).
\end{equation}
More generally, for every fixed $a>0$, the same rate holds for any cap
$Q$ satisfying
$L_T+1\le Q\le[\log(edT\kappa)]^a$, whenever this interval is nonempty.
Constants implicit in the suppressed logarithmic factors may depend on $a$.
\end{corollary}

\Cref{tab:phases} states the resulting regimes in terms of the total regret.
It separates the transition in the geometric term from the point at which
that term actually exceeds the online-learning cost.

\begin{table}[t]
\centering
\caption{Minimax expected regret with a strict polylogarithmic separation
budget. Here $M=\min\{T,d\}$; rates suppress logarithmic factors and universal
constants. Boundary cases agree up to constants.}
\label{tab:phases}
\small
\begin{tabularx}{\textwidth}{@{}P{0.22\textwidth}Y P{0.24\textwidth}@{}}
\toprule
Size of $M$ & Asphericity range & Regret \\
\midrule
$M\le\sqrt T$ & Every $\kappa\ge1$ & $\wtTheta(GR\sqrt T)$ \\
\addlinespace
\(M>\sqrt T\) & $1\le\kappa\le\sqrt{T/M}$ & $\wtTheta(GR\sqrt T)$ \\
\addlinespace
\(M>\sqrt T\) & $\sqrt{T/M}<\kappa<\sqrt M$ & $\wtTheta(GR\kappa\sqrt M)$ \\
\addlinespace
\(M>\sqrt T\) & $\kappa\ge\sqrt M$ & $\wtTheta(GRM)$ \\
\bottomrule
\end{tabularx}
\end{table}

\section{Comparison with prior work}\label{sec:related}

Projection-based online gradient descent supplies the $GR\sqrt T$ benchmark,
but computing a projection is not a unit-cost operation in our model
\cite{zinkevich2003}. Oracle-based alternatives must therefore be compared
using both their regret and the information required to produce an action.
\Cref{tab:related} uses a common normalization and distinguishes a per-round
cap from a total budget.

\begin{table}[t]
\centering
\caption{Results for fixed-domain full-information OCO under a known ball
sandwich. Regret is divided by $GR$; logarithms are suppressed. The
membership rows are distinct algorithms. ``Per round'' refers to a bound
on every round; ``total'' permits uneven allocation.
Here $H_M:=\min\{M,\kappa\sqrt M\}$ and $M=\min\{T,d\}$.}
\label{tab:related}
\small
\renewcommand{\arraystretch}{1.18}
\begin{tabularx}{\textwidth}{@{}P{0.28\textwidth}P{0.15\textwidth}P{0.24\textwidth}Y@{}}
\toprule
Work & Oracle & Calls & Regret divided by $GR$ \\
\midrule
Mhammedi (2022), Thm.~8 & Membership & $\wtO(d)$ per round & $\wtO(\kappa\sqrt T)$ \\
Mhammedi (2022), Thm.~15 & Membership & $\wtO(1)$ per round & $\wtO(\kappa\sqrt{dT})$ \\
Gauge reduction\newline Mhammedi (2022, 2025) & Separation & $\wtO(1)$ per round & $O(\kappa\sqrt T)$ \\
Garber--Kretzu (2022), Thm.~14 & Separation & $O(T)$ total & $O(\kappa\sqrt T)$ \\
Mhammedi (COLT 2025), Thm.~4.1 & Separation & $\wtO(1)$ per round & $\wtO(\sqrt{dT}+\kappa d)$ \\
Mhammedi (NeurIPS 2025), Thm.~4.1 & Separation & $\wtO(1)$ per round & $\wtO(\sqrt{dT}+d^2)$ \\
\midrule
This work, \Cref{thm:upper} & Separation & $1$ per round & $\wtO(\min\{T,\sqrt T+d\})$ \\
This work, \Cref{cor:main} & Separation & $Q_\star=O(\log(e\kappa T))$ per round &
$\begin{gathered}\wtTheta(\min\{T,\\[-2pt]\sqrt T+H_M\})\end{gathered}$ \\
\bottomrule
\end{tabularx}
\end{table}

\paragraph{Membership and separation reductions.}
Mhammedi~\cite{mhammedi2022membership} reduces OCO on a general body to OCO
on a ball using approximate gauge computations. The dimension-independent
$\kappa\sqrt T$ regret bound in Theorem~8 uses $\wtO(d)$ membership calls
per round; Theorem~15 reduces the calls to $\wtO(1)$ and pays a $\sqrt d$
factor in expected regret. With separation access, the gauge reduction
attains $O(\kappa\sqrt T)$ using only logarithmically many calls. A direct
deterministic analysis is given in Mhammedi~\cite[Lemma~4.1 and
Remark~4.1]{mhammedi2025sep}. Our fallback is a self-contained version of
this established reduction, adjusted to an unknown number of selected rounds.

\paragraph{Total-budget methods.}
Garber and Kretzu~\cite[Theorem~14]{garberkretzu2022} obtain a separation-based
algorithm with adaptive-regret guarantees and $O(T)$ total oracle calls.
Restoring the dependence on the ball sandwich gives
$O(\kappa GR\sqrt T)$ regret. For the displayed total-call bound, their
parameter choice applies when $T>16\kappa^2$; playing the origin handles
the remaining horizons at the same regret order. This total budget is a
different restriction from the hard cap considered here. Their guarantees
over all time intervals also address a stronger comparator criterion than
our static regret.

\paragraph{Strict-budget separation methods.}
Mhammedi~\cite[Theorem~4.1]{mhammedi2025sep} proves
$\wtO(GR[\sqrt{dT}+\kappa d])$ regret with a polylogarithmic per-round
separation budget. The subsequent ellipsoid
method~\cite[Theorem~4.1]{mhammedi2025ellipsoid} obtains
$\wtO(GR[\sqrt{dT}+d^2])$, with only logarithmic dependence on $\kappa$.
Our bound replaces the dimension-dependent leading term by $\sqrt T$ and
identifies the additional geometric cost up to logarithms. This comparison
concerns oracle complexity. The ellipsoid method also provides efficient
per-round matrix computations, which the centroid implementation here does
not match.

The first membership row is a high-probability bound, while the second is an
expected bound for each comparator; these distinctions do not turn either
row into a dimension-independent, polylogarithmic membership guarantee.
The separation reductions and the two cited separation algorithms are
deterministic. Our localization upper bound is in expectation, and its
high-probability event gives a guarantee simultaneously for all comparators.

\paragraph{Other oracle-budget tradeoffs.}
Aggarwal~\cite{aggarwal2026} studies total and per-round query budgets for
exact linear optimization access. That oracle returns a feasible optimizer
of a supplied linear objective. Its information interface differs from a
separating hyperplane at a supplied point, so those bounds address a
different oracle model from the one in \cref{def:sep}.

\section{A deterministic feasible fallback}\label{sec:backup}

The hybrid algorithm needs a learner that can be paused between selected
rounds. Its guarantee must hold for every realized number of selected rounds,
because that number is determined by localization. We now construct such a
learner using only exact separation.

\subsection{Radial feasibility and a surrogate gradient}

Let $y\in RB_2^d$ be an auxiliary action. If a separation query certifies
$y\in\K$, play $x=y$ and set $h=0$. Otherwise initialize $a=1$, $b=\kappa$,
and retain the separator at $y/a$. The point $y/b$ is feasible because it
lies in $rB_2^d$. Repeatedly bisect $[a,b]$: query $y/c$ at
$c=(a+b)/2$, replace $b$ by $c$ after a feasible response, and replace $a$
by $c$ and retain the new separator after an infeasible response. Stop when
$b-a\le1/T$. At most $L_T$ calls, including the first query, are needed.
The case $\kappa=1$ has a known feasible ball and needs no bisection.

Let $n$ be the unit normal at the final infeasible point $y/a$, and define
\begin{equation}\label{eq:polar}
 x:=y/b,
 \qquad h:=\frac{n}{\ip{n}{y/a}}.
\end{equation}
The separator and the inner ball imply
\begin{equation}\label{eq:polar-properties}
 x\in\K,\qquad
 \norm{h}_2\le1/r,\qquad
 \ip{h}{u}\le1\quad(u\in\K),\qquad
 \ip{h}{y}=a.
\end{equation}
Indeed, $\ip{n}{y/a}>\sup_{u\in\K}\ip{n}{u}\ge r$.
When the loss subgradient $g$ is revealed, set
\begin{equation}\label{eq:surrogate}
 \widetilde g:=g+(-\ip{g}{x})_+h,
 \qquad (z)_+:=\max\{z,0\}.
\end{equation}

\begin{lemma}[One-step reduction]\label{lem:radial}
For every $u\in\K$ and every $g$ with $\norm{g}_2\le G$,
\begin{equation}\label{eq:radial}
 \ip{g}{x-u}\le\ip{\widetilde g}{y-u}+GR/T,
 \qquad
 \norm{\widetilde g}_2\le(1+\kappa)G.
\end{equation}
The action $x$ is feasible regardless of the subsequent loss.
\end{lemma}

\begin{proof}
The feasible case has $x=y$ and $h=0$, so the claim is immediate.
In the remaining case, $y=bx$ with $b\ge1$. If $\ip{g}{x}\ge0$, then
$\ip{g}{x-y}=(1-b)\ip{g}{x}\le0$. Otherwise put
$c=-\ip{g}{x}>0$. By \eqref{eq:polar-properties},
\begin{align*}
 \ip{g}{x-u}-\ip{g+ch}{y-u}
 &=c(b-1)-c(a-\ip{h}{u})\\
 &\le c(b-a)\le GR/T.
\end{align*}
Finally $c\le GR$, so
$\norm{\widetilde g}_2\le G+GR/r=(1+\kappa)G$.
\end{proof}

\subsection{Anytime guarantee on selected rounds}

Set $H=(1+\kappa)G$. On virtual round $j$, the fallback maintains
$y_j\in RB_2^d$, computes its feasible action by the preceding procedure,
and, after receiving $g_j$, performs
\begin{equation}\label{eq:ogd}
 y_{j+1}=\proj_{RB_2^d}(y_j-\eta_j\widetilde g_j),
 \qquad \eta_j=\frac{2R}{H\sqrt j},\qquad y_1=0.
\end{equation}
This projection is onto an explicit ball and needs no oracle access.

\begin{lemma}[Fallback guarantee]\label{lem:backup}
The deterministic learner $\Back$ uses at most $L_T$ separation queries
per virtual action. For every $0\le F\le T$, every realized sequence of
subgradients, and every $u\in\K$,
\begin{equation}\label{eq:backup}
 \sum_{j=1}^F\ip{g_j}{x_j-u}\le7\kappa GR\sqrt F.
\end{equation}
\end{lemma}

\begin{proof}
Nonexpansiveness of projection gives
\[
 \ip{\widetilde g_j}{y_j-u}
 \le\frac{\norm{y_j-u}_2^2-\norm{y_{j+1}-u}_2^2}{2\eta_j}
       +\frac{\eta_jH^2}{2}.
\]
The step sizes are nonincreasing and $\norm{y_j-u}_2\le2R$.
Summation therefore bounds the distance terms by $2R^2/\eta_F$.
Since $\sum_{j=1}^Fj^{-1/2}\le2\sqrt F$,
\[
 \sum_{j=1}^F\ip{\widetilde g_j}{y_j-u}
 \le RH\sqrt F+2RH\sqrt F=3RH\sqrt F.
\]
By \cref{lem:radial}, the original regret adds at most $GRF/T$.
For $F\ge1$, use $F/T\le\sqrt F$ and $1+\kappa\le2\kappa$.
The claim for $F=0$ is immediate.
\end{proof}

The horizon $T$ is used only to fix the bisection accuracy and its query cap;
the step sizes do not require advance knowledge of $F$. Because
\eqref{eq:backup} is pathwise and holds for every prefix, it applies without
an optional-stopping argument when localization determines which rounds
are sent to the fallback.

\section{Gaussian localization and the upper bound}\label{sec:upper}

\subsection{Algorithm and query allocation}

The primary learner stores an explicit convex outer approximation $D$ and
the sum $s$ of gradients on previously accepted rounds. Initially
$D=RB_2^d$ and $s=0$. For parameters $\lambda,\mu>0$, define
\begin{equation}\label{eq:gibbs}
 p_{D,s}(x):=\frac{1}{Z(D,s)}
 \exp\!\left(-\frac{\lambda\norm{x}_2^2/2+\ip{s}{x}}{\mu}\right)
 \ind_D(x),
 \qquad m(D,s):=\E_{p_{D,s}}X.
\end{equation}
Each localization query uses an approximate centroid $\widehat m$ of this
density. A feasible answer supplies the action. An infeasible answer adds
the valid halfspace through $\widehat m$ to $D$ and the learner tries again.
If all $q$ queries fail, the fallback supplies the action.

\begin{figure}[t]
\hrule\smallskip
\refstepcounter{algorithm}\label{alg:hybrid}
\noindent\textbf{Algorithm \thealgorithm. Gaussian localization with a feasible fallback}
\smallskip\hrule\medskip
\small
\noindent\textbf{Input:} $Q\ge L_T+1$, $q=Q-L_T$, $\lambda,\mu>0$.
\par\smallskip
\noindent\textbf{Initialize:} $D=RB_2^d$, $s=0$, and the learner $\Back$.
\par\smallskip
\noindent\textbf{On each round $t=1,\ldots,T$:}
\begin{enumerate}[label=\arabic*.,leftmargin=1.7em,itemsep=3pt,topsep=4pt]
 \item Set $\mathsf{accepted}=\mathsf{false}$.
 \item For at most $q$ attempts, compute an approximate centroid $\widehat m$
 of $p_{D,s}$ and query $\Sep_\K(\widehat m)$:
 \begin{itemize}[leftmargin=1.4em,itemsep=2pt,topsep=2pt]
  \item If the answer is \texttt{inside}, set $w_t=\widehat m$ and
  $\mathsf{accepted}=\mathsf{true}$, and stop the attempts.
  \item Otherwise, for the returned normal $a$, replace $D$ by
  $D\cap\{x:\ip{a}{x}\le\ip{a}{\widehat m}\}$.
 \end{itemize}
 \item If all attempts fail, obtain the next feasible action $w_t$ of
 $\Back$, using at most $L_T$ additional separation calls.
 \item Play $w_t$ and observe $g_t$.
 \item On an accepted round, set $s\gets s+g_t$. Otherwise, update $\Back$
 with $g_t$ on its next virtual round.
\end{enumerate}
\smallskip\hrule
\end{figure}

Every real round uses at most $q+L_T=Q$ queries. The fallback is invoked
only after the primary queries fail and before the current loss is known.
Its internal state depends only on the gradients from earlier failure
rounds. This gives precisely the virtual sequence covered by
\cref{lem:backup}. For the origin-fallback variant, omit $\Back$, allow
$Q$ primary queries, and play $0$ whenever all of them fail.

At every step, $D$ contains $\K$ and has the explicit form
\[
 D=RB_2^d\cap\bigcap_{j\le C}\{x:\ip{a_j}{x}\le b_j\}.
\]
Testing membership in $D$ requires only the stored inequalities.
In particular, centroid sampling does not query membership in the unknown
body $\K$. Exact feasibility of the played actions follows from either an
\texttt{inside} response or the deterministic fallback construction, even
when a centroid estimate is inaccurate.

\subsection{Dimension-independent stability}

Define $\Phi(D,s):=\mu\log Z(D,s)$.

\begin{lemma}[Gaussian stability]\label{lem:stability}
For every convex $D$ with nonempty interior and every $s,g\in\R^d$,
\begin{equation}\label{eq:stability}
 \Phi(D,s+g)-\Phi(D,s)
 \le-\ip{g}{m(D,s)}+\frac{\norm{g}_2^2}{2\lambda}.
\end{equation}
Equivalently, the covariance of $p_{D,s}$ satisfies
$\Cov(X)\preceq(\mu/\lambda)I$.
\end{lemma}

\begin{proof}
Completing the square gives
\[
 Z(D,s)=e^{\norm{s}_2^2/(2\lambda\mu)}
 \int_D\exp\!\left(-\frac{\lambda}{2\mu}
                  \norm{x+s/\lambda}_2^2\right)\,dx.
\]
The integral is log-concave as a function of $s$, by preservation of
log-concavity under marginalization~\cite[Theorem~5.1]{lovaszvempala2007}.
Consequently $\nabla_s^2\log Z\preceq(\lambda\mu)^{-1}I$.
Differentiation under the integral gives
$\nabla_s\log Z=-m(D,s)/\mu$ and
$\nabla_s^2\log Z=\Cov(X)/\mu^2$.
The smoothness inequality proves \eqref{eq:stability}.
\end{proof}

The covariance bound is in operator norm. Replacing it by a trace bound
would introduce an unnecessary factor of $d$ into the stability cost.

\subsection{Constant-depth cuts from approximate centroids}

Let $m=m(D,s)$ and $\Sigma=\Cov_{p_{D,s}}(X)$. The centroid procedure is
chosen so that, simultaneously for every primary query,
\begin{equation}\label{eq:centroid-accuracy}
 \norm{\widehat m-m}_2\le R/T,
 \qquad
 \norm{\Sigma^{-1/2}(\widehat m-m)}_2\le\frac1{2e}.
\end{equation}
\Cref{app:sampling} constructs such a procedure with overall failure
probability at most $1/(4T^2)$. The two inequalities have distinct roles:
Euclidean error controls regret, and Mahalanobis error controls cut depth.

\begin{lemma}[Partition decrease]\label{lem:cut}
If \eqref{eq:centroid-accuracy} holds and the oracle rejects $\widehat m$,
then the updated outer body $D'$ satisfies
\begin{equation}\label{eq:cut}
 \Phi(D',s)-\Phi(D,s)\le-\gamma\mu,
 \qquad \gamma:=-\log(1-1/(2e))>0.
\end{equation}
\end{lemma}

\begin{proof}
Let $a$ be the returned unit normal and put
$\sigma_a=(a^\top\Sigma a)^{1/2}>0$.
The random variable $Y=\ip{a}{X-m}/\sigma_a$ is log-concave, has mean zero,
and has variance one. A halfline through its mean has probability at least
$1/e$, and its density is bounded above by one
\cite[Lemmas~5.4 and~5.5]{lovaszvempala2007}.
By Cauchy--Schwarz,
\[
 \frac{|\ip{a}{\widehat m-m}|}{\sigma_a}
 \le\norm{\Sigma^{-1/2}(\widehat m-m)}_2\le1/(2e).
\]
Moving the threshold from $m$ to $\widehat m$ therefore leaves probability
at least $1/(2e)$ in each open halfspace. The valid cut retains at most
$1-1/(2e)$ of the probability mass, proving \eqref{eq:cut}.
\end{proof}

\subsection{A potential bound for the accepted rounds}

Let $\mathcal A$ and $\mathcal F$ be the accepted and failure rounds,
respectively. Write $F=|\mathcal F|$ and let $C$ be the number of rejected
primary queries. In the hybrid, $C\ge qF$.

\begin{lemma}[Accepted-round regret]\label{lem:accepted}
Set $\lambda=G\sqrt T/R$. On the event that all centroid estimates satisfy
\eqref{eq:centroid-accuracy}, simultaneously for every $u\in\K$,
\begin{equation}\label{eq:accepted}
 \sum_{t\in\mathcal A}\ip{g_t}{w_t-u}
 \le GR(\sqrt T+3)+\mu D_T-\gamma\mu C.
\end{equation}
\end{lemma}

\begin{proof}
Apply \cref{lem:stability} to each accepted gradient update and
\cref{lem:cut} to each rejected query. Summing in the actual chronological
order, including the cuts between successive gradient updates, gives
\begin{align}\label{eq:telescoping}
 \sum_{t\in\mathcal A}\ip{g_t}{m_t-u}
 \le{}&\Phi(D_0,0)-\Phi(D_{\rm fin},s_{\rm fin})
       -\ip{s_{\rm fin}}{u}\\
 &+\frac{TG^2}{2\lambda}-\gamma\mu C.\nonumber
\end{align}
Here $m_t$ is the centroid at the accepted query, and
$s_{\rm fin}=\sum_{t\in\mathcal A}g_t$.
Replacing $m_t$ by $w_t$ adds at most $GR$.

To bound the remaining partition terms, put $\delta=1/T$. Convexity and
\eqref{eq:sandwich} imply
\[
 U_u:=(1-\delta)u+\delta rB_2^d\subseteq\K\subseteq D_{\rm fin}.
\]
For $x\in U_u$, $\norm{x}_2\le R$ and $\norm{x-u}_2\le2R/T$.
Since $\norm{s_{\rm fin}}_2\le TG$, integration over $U_u$ yields
\[
 Z(D_{\rm fin},s_{\rm fin})
 \ge T^{-d}\vol(rB_2^d)
 \exp\!\left(-\frac{\lambda R^2/2+\ip{s_{\rm fin}}{u}+2GR}{\mu}\right).
\]
Also $Z(D_0,0)\le\vol(RB_2^d)$. Therefore
\[
 \Phi(D_0,0)-\Phi(D_{\rm fin},s_{\rm fin})-\ip{s_{\rm fin}}{u}
 \le\frac{\lambda R^2}{2}+2GR+\mu d\log(\kappa T).
\]
Substitute in \eqref{eq:telescoping}. The choice of $\lambda$ makes
$\lambda R^2/2+TG^2/(2\lambda)=GR\sqrt T$.
Finally, $d\log(\kappa T)\le D_T$.
\end{proof}

\subsection{Paying for failed localization}

\begin{proof}[Proof of \cref{thm:upper}]
For the origin-fallback strategy, use all $Q$ primary queries per round.
The failure rounds contribute at most $GRF$, and $C\ge QF$.
Choosing $\mu=GR/(\gamma Q)$ in \eqref{eq:accepted} gives, on the good
centroid event,
\begin{equation}\label{eq:origin-explicit}
 \Reg_T\le GR(\sqrt T+3)+\frac{GRD_T}{\gamma Q}.
\end{equation}

For the hybrid, put $q=Q-L_T$ and use \cref{lem:backup}. Uniformly over
$u\in\K$, the combined regret is at most
\[
 GR(\sqrt T+3)+\mu D_T-\gamma\mu qF+7\kappa GR\sqrt F.
\]
For $A,B>0$, $A\sqrt F-BF\le A^2/(4B)$. Hence
\[
 \Reg_T\le GR(\sqrt T+3)+\mu D_T
       +\frac{49\kappa^2G^2R^2}{4\gamma\mu q}.
\]
The choice
$\mu=7\kappa GR/(2\sqrt{\gamma qD_T})$ gives
\begin{equation}\label{eq:hybrid-explicit}
 \Reg_T\le GR(\sqrt T+3)
       +7\kappa GR\sqrt{\frac{D_T}{\gamma q}}.
\end{equation}
This use of the fallback does not require a deterministic value of $F$:
its bound holds on every realized virtual sequence.

All played actions remain feasible if centroid estimation fails. By
\eqref{eq:regret}, regret is then at most $2GRT$. The failure probability
$1/(4T^2)$ adds at most $GR/(2T)$ in expectation.
Running $\Back$ alone gives $7\kappa GR\sqrt T$ using at most $L_T$ calls
per round. Playing $0$ throughout gives $GRT$ with no queries.
Selecting the best guaranteed strategy proves both claims.
\end{proof}

\section{Hidden facets and the lower bound}\label{sec:lower}

The lower bound isolates the cost of discovering feasible actions. The loss
does not change with time; all uncertainty is in the body and its oracle.

\subsection{A family with tunable comparator advantage}

Fix $\kappa>1$ and $1\le m\le d-1$. Write points as
$z=(a,x,y)\in\R\times\R^m\times\R^{d-m-1}$ and let
$V=(V_1,\ldots,V_m)$ be uniform on $\{-1,+1\}^m$. Define
\begin{equation}\label{eq:facet-parameters}
 \alpha=\sqrt{\frac{2}{\kappa^2-1}},\qquad
 \beta=r\sqrt{1+\alpha^2},\qquad
 \tau=\frac{\beta}{\alpha}
 =R\sqrt{\frac{\kappa^2+1}{2\kappa^2}}.
\end{equation}
The hidden-facet body is
\begin{equation}\label{eq:body}
 \K_V^{(m)}:=RB_2^d\cap
 \bigcap_{i=1}^m\{(a,x,y):\alpha a-V_ix_i\le\beta\}.
\end{equation}

\begin{lemma}[Geometry of the hidden-facet family]\label{lem:geometry}
Each body in \eqref{eq:body} satisfies \eqref{eq:sandwich}. An optimizer
of $\max_{z\in\K_V^{(m)}}a$ is
\begin{equation}\label{eq:comparator}
 u_V=(\tau+\Delta_{m,\kappa},\alpha\Delta_{m,\kappa}V,0),
 \qquad \Delta_{m,\kappa}=R\delta_{m,\kappa}.
\end{equation}
Moreover,
\begin{equation}\label{eq:delta-lower}
 \delta_{m,\kappa}\ge\frac18
 \min\left\{1-\kappa^{-2},\frac{\kappa^2-1}{\kappa\sqrt m}\right\}.
\end{equation}
\end{lemma}

\begin{proof}
Each facet normal has norm $\sqrt{1+\alpha^2}$ and offset
$r\sqrt{1+\alpha^2}$, so its halfspace contains $rB_2^d$.
The outer inclusion follows from intersection with $RB_2^d$.

For fixed $a>\tau$, the feasible vector of smallest norm has
$x_i=V_i\alpha(a-\tau)$ and $y=0$. Thus the largest feasible $a$ is the
root above $\tau$ of
\[
 a^2+m\alpha^2(a-\tau)^2=R^2.
\]
Let $\rho=\tau/R$. The positive root for $\delta=(a-\tau)/R$ is
\[
 \delta=\frac{1-\rho^2}
 {\rho+\sqrt{\rho^2+(1+m\alpha^2)(1-\rho^2)}}.
\]
Using $1-\rho^2=(1-\kappa^{-2})/2$ and
$\alpha^2(1-\rho^2)=\kappa^{-2}$ gives \eqref{eq:delta}.
Finally,
\[
 \delta_{m,\kappa}
 \ge\frac{(1-\kappa^{-2})/2}{2+\sqrt m/\kappa}.
\]
If $\sqrt m/\kappa\le2$, the denominator is at most $4$;
otherwise it is at most $2\sqrt m/\kappa$. These two cases imply
\eqref{eq:delta-lower}.
\end{proof}

Choose the same loss on every round:
\begin{equation}\label{eq:fixed-loss}
 f_t(a,x,y)=-Ga.
\end{equation}
Every feasible action with $a\le\tau$ then loses at least
$G\Delta_{m,\kappa}$ relative to $u_V$. Since $u_V$ maximizes $a$, regret
on other rounds is nonnegative as well.

\subsection{A fixed oracle and its conditional information}

For each $V$, fix the following oracle once and for all. Given
$z=(a,x,y)$, it first checks the outer ball. If $\norm{z}_2>R$, it returns
$z/\norm{z}_2$. Otherwise it returns the normalized normal
$(\alpha e_0-V_ie_i)/\sqrt{1+\alpha^2}$ for the least index $i$ with
$\alpha a-V_ix_i>\beta$. If no such index exists, it returns
\texttt{inside}. Each response is a valid exact strong separator or a
correct feasibility certificate. In particular, the construction does not
choose an oracle retrospectively from the learner's future actions.

\begin{lemma}[Posterior subcube]\label{lem:subcube}
Fix a deterministic learner. Conditional on any positive-probability
adaptive transcript, $V$ is uniform on
$S_1\times\cdots\times S_m$, where each
$S_i$ is one of $\{-1\}$, $\{+1\}$, or $\{-1,+1\}$.
\end{lemma}

\begin{proof}
The initial prior has this form. Given a transcript, the next query is
fixed. A ball response imposes no restriction on $V$. A response from facet
$j$ requires every earlier facet to be satisfied, fixes $V_j$ through the
returned normal, and imposes no condition on later coordinates. An
\texttt{inside} response requires each facet to be satisfied. Each
condition concerns a single sign. Its intersection with a coordinate
subcube is another such subcube, and all surviving sign vectors retain
equal probability. The losses in \eqref{eq:fixed-loss} reveal no additional
information. Induction proves the claim.
\end{proof}

\begin{lemma}[Information per query]\label{lem:information}
Let $J$ count previously unfixed signs that become fixed after a single
oracle response. Conditional on any previous transcript,
\begin{equation}\label{eq:information}
 \Prob(J\ge k\mid\text{transcript})\le2^{-(k-1)}
 \quad(k\ge1),\qquad
 \E[J\mid\text{transcript}]\le2.
\end{equation}
\end{lemma}

\begin{proof}
If the query is outside the outer ball, $J=0$. Otherwise inspect the
still-unknown coordinates in oracle order.
At each coordinate either both signs satisfy its facet, exactly one sign
does, or neither sign does. A coordinate of the first type reveals
nothing. At a coordinate of the second type, continuing past that facet
requires its independent uniform sign to equal the satisfying value, with
conditional probability $1/2$. A coordinate of the third type stops the
query and can reveal only the final sign. Already fixed coordinates can
only stop the query earlier.
Thus revealing at least $k$ new signs requires passing at least $k-1$
independent tests of probability $1/2$. This proves the tail bound.
Summing $\E J=\sum_{k\ge1}\Prob(J\ge k)$ gives the second claim.
\end{proof}

\subsection{Why feasibility limits an uninformed action}

If sign $V_i$ is not fixed by a transcript, both of its values remain
possible. Any action feasible for every consistent body must satisfy
\[
 \alpha a-x_i\le\beta,
 \qquad \alpha a+x_i\le\beta.
\]
Adding these inequalities gives $a\le\tau$.

This statement also applies to randomized learners with almost-sure
feasibility. For each of the finitely many bodies, the set of internal
random seeds on which the learner remains feasible has probability one.
Intersect these sets over all $2^m$ bodies and fix a seed in the
intersection. The learner is then deterministic and feasible for every
consistent body. The same argument can be averaged over the seed after
applying \cref{lem:information}.

\begin{proof}[Proof of \eqref{eq:lower-exact}]
By the time the action on round $t$ is selected, at most $Qt$ oracle
responses have been received. Pad a shorter query sequence with queries
at the known feasible origin if necessary. By
\cref{lem:information}, the expected number of fixed signs is at most
$2Qt$. Therefore
\[
 \Prob(\text{all $m$ signs are fixed before action $t$})\le\frac{2Qt}{m}.
\]
For $t\le m/(4Q)$, with probability at least $1/2$ an unknown sign remains
and feasibility forces $a_t\le\tau$. Summing the resulting regret under
\eqref{eq:fixed-loss} gives
\[
 \E\Reg_T\ge\frac{GR\delta_{m,\kappa}}2
 \left\lfloor\min\left\{T,\frac{m}{4Q}\right\}\right\rfloor.
\]
This holds under the uniform prior for every randomized admissible learner,
so some fixed body has at least this expected regret.

For completeness, take the known body $RB_2^d$ and independent Rademacher
signs $\epsilon_1,\ldots,\epsilon_T$. Use losses
$f_t(w)=G\epsilon_t w_1$, with the current sign unrevealed before $w_t$.
The learner's expected loss is zero, while the best fixed comparator has
loss $-GR|S_T|$, where $S_T=\sum_t\epsilon_t$.
Since $\E S_T^2=T$ and $\E S_T^4=3T^2-2T\le3T^2$, the
Paley--Zygmund inequality gives
$\Prob(S_T^2\ge T/2)\ge1/12$. Hence
$\E|S_T|\ge\sqrt T/(12\sqrt2)$.
The maximum of this ball lower bound and the hidden-facet bounds proves
\eqref{eq:lower-exact}. For $\kappa=1$, only the ball argument is needed.
\end{proof}

\subsection{The simplified query frontier}

\begin{proof}[Proof of \eqref{eq:lower-simple}]
Suppose $\kappa\ge2$ and choose the integer
$m=\min\{d-1,QT\}$. Then $m/Q=N_Q$.
If $N_Q<8$, the proposed geometric term is at most $8$ and is absorbed,
up to a universal constant, by the ball lower bound.
If $N_Q\ge8$, then
\[
 \left\lfloor\min\left\{T,\frac{m}{4Q}\right\}\right\rfloor
 =\lfloor N_Q/4\rfloor\ge N_Q/8.
\]
Moreover, \cref{lem:geometry} implies
$\delta_{m,\kappa}\ge c\min\{1,\kappa/\sqrt m\}$.
Thus the hidden-facet contribution is at least
\[
 cGR\min\left\{N_Q,\kappa\frac{N_Q}{\sqrt m}\right\}
 =cGR\min\left\{N_Q,\kappa\sqrt{\frac{N_Q}{Q}}\right\}.
\]
Combine this with the ball lower bound using
$\max\{a,b\}\ge(a+b)/2$, and take the outer minimum with $T$.
\end{proof}

\section{Consequences and the unresolved budget dependence}\label{sec:consequences}

\begin{proof}[Proof of \cref{cor:main}]
First consider the upper bound. The reserve condition gives $q\ge1$.
If $d\le T$, the two localization branches in
\eqref{eq:upper-hybrid} give
$\wtO(GR[\sqrt T+\min\{d,\kappa\sqrt d\}])$, and the origin strategy
caps this at $GRT$. If $d>T$, the minimum of the fallback-only and origin
strategies is $O(GR\min\{\kappa\sqrt T,T\})$, which is within a constant
factor of the asserted expression with $M=T$.

For the lower bound, $d=1$ or $1\le\kappa<2$ is settled by the known-ball
bound because $\min\{M,\kappa\sqrt M\}\le2\sqrt T$ in these cases.
For $d\ge2$ and $\kappa\ge2$, observe that
\[
 N_Q=\min\{T,(d-1)/Q\}\ge\frac{M}{2Q}.
\]
Consequently
\[
 \min\left\{N_Q,\kappa\sqrt{N_Q/Q}\right\}
 \ge\frac1{2Q}\min\{M,\kappa\sqrt M\}.
\]
Apply \eqref{eq:lower-simple}. A factor $Q$ bounded by a fixed power of
$\log(edT\kappa)$ is absorbed by the tilde notation. This proves the
claim for the stated budget range and, in particular, for $Q_\star$.
\end{proof}

\paragraph{Larger query budgets.}
The dependence on $Q$ should not be suppressed when $Q$ grows faster than
a polylogarithm. For example, when $d\le QT$ and the intermediate
asphericity terms are relevant, the upper and lower geometric terms are,
respectively,
\begin{equation}\label{eq:Qgap}
 \wtO\!\left(GR\kappa\sqrt{\frac d{Q-L_T}}\right),
 \qquad
 \Omega\!\left(GR\kappa\frac{\sqrt{d-1}}{Q}\right).
\end{equation}
Their ratio, apart from logarithms and constants, is
$Q/\sqrt{Q-L_T}$. If $Q\ge2L_T$, this is $O(\sqrt Q)$.
When $d>QT$, the fallback-only term $\kappa\sqrt T$ and the corresponding
lower term $\kappa\sqrt{T/Q}$ lead to the same $\sqrt Q$ comparison.
Near the threshold $Q=L_T+1$, however, one must retain the reserve
explicitly; a uniform $\sqrt Q$ statement would be inaccurate.
Closing this budget-dependent gap remains open.

\paragraph{One call per round.}
The origin-fallback algorithm is valid with $Q=1$ and achieves
$\wtO(GR[\sqrt T+d])$, capped by $GRT$.
Its guarantee does not establish the hybrid asphericity improvement at
that literal cap. The fallback queries cannot be omitted from the
accounting, and \cref{cor:main} makes no such claim.

\paragraph{Other access models.}
An exact separation call simulates an exact membership call by discarding
the normal. Thus our separation lower bound also applies to learners with
only membership access under the same cap. The upper bound does not
transfer: both the primary cuts and the low-query fallback use separating
normals. A weak oracle that may certify points at distance at most
$\eps$ from $\K$ changes the feasibility requirement. Simply accepting
its certificates can produce infeasible actions. A theorem for that model
must specify the allowed constraint violation and the separation margin.

\section{Discussion}

The analysis identifies two distinct sources of regret under limited
separation access. Online stability costs $GR\sqrt T$ even on a known
ball. Geometric uncertainty contributes an additional term that depends on
both the number of coordinates that can remain hidden and the comparator
advantage permitted by the ball sandwich. Gaussian localization and the
hidden-facet family meet at this scale when the per-round budget is
polylogarithmic.

Two questions remain particularly direct. The first is the exact dependence
on larger budgets, where \eqref{eq:Qgap} records a gap rather than a minimax
characterization. The second is computational: can one attain the same
regret with the per-round arithmetic efficiency of ellipsoid-based
methods? The elementary sampler below establishes oracle implementability,
but not an efficient numerical algorithm. A finite-precision result would
also require quantitative assumptions on separator encoding and numerical
certification. These issues do not affect the strict unknown-body query
count proved here, but they matter for turning the minimax construction
into a practical method.

\appendix
\section{Centroid estimation without unknown-body queries}\label{app:sampling}

This appendix makes the implementation used by \cref{thm:upper} explicit.
Only the ball and the stored halfspaces defining $D$ are consulted.
The construction proves finite expected internal work, not a polynomial
arithmetic bound. More sophisticated log-concave sampling algorithms
\cite{lovaszvempala2007} may be substituted, provided their approximation
errors are controlled as described below.

\subsection{Exact samples by rejection}

For a current state $(D,s)$, let
$V(x)=\lambda\norm{x}_2^2/2+\ip{s}{x}$.
Draw $X$ uniformly from $RB_2^d$ and $U$ uniformly from $[0,1]$,
independently. Accept $X$ if $X\in D$ and
\begin{equation}\label{eq:rejection}
 U\le\exp\!\left(-\frac{V(X)+R\norm{s}_2}{\mu}\right).
\end{equation}
The right-hand side lies in $[0,1]$ on $RB_2^d$ because
$V(X)\ge-R\norm{s}_2$. The density of an accepted point is exactly
$p_{D,s}$. Since $D$ contains $rB_2^d$, the acceptance probability is at least
\begin{equation}\label{eq:accept-prob}
 p_{\min}(s):=
 \kappa^{-d}\exp\!\left(-\frac{\lambda R^2/2+2R\norm{s}_2}{\mu}\right)>0.
\end{equation}
Restarting with independent randomness therefore produces independent
exact samples with finite expected work. Also $\norm{s}_2\le TG$, so a
known uniform positive lower bound is obtained by replacing
$\norm{s}_2$ with $TG$ in \eqref{eq:accept-prob}.

If a deterministically bounded number of arithmetic trials is desired,
stop each sample after $K$ proposals and output $0$ on timeout. Its law
then differs from the target by at most $(1-p_{\min})^K$ in total
variation. Choosing $K$ using the known uniform lower bound makes this
error as small as required. This bound can be exponential, which is why
no polynomial-time assertion is made.

\subsection{Two norms and a uniform success event}

Let $X_1,\ldots,X_N$ be independent exact samples from $p_{D,s}$ and
$\overline X=N^{-1}\sum_iX_i$. With centroid $m$ and covariance $\Sigma$,
\begin{equation}\label{eq:sample-moments}
 \E\norm{\overline X-m}_2^2
 =\frac{\operatorname{tr}\Sigma}{N}\le\frac{R^2}{N},
 \qquad
 \E\norm{\Sigma^{-1/2}(\overline X-m)}_2^2=\frac dN.
\end{equation}
Positive definiteness follows because the density is positive on a body
containing a nonempty ball. Markov's inequality gives
\[
 \Prob\!\left(\norm{\overline X-m}_2>R/T
 \text{ or }
 \norm{\Sigma^{-1/2}(\overline X-m)}_2>1/(2e)\right)
 \le\frac{T^2+4e^2d}{N}.
\]
Let $J_{\max}=QT$ bound the total number of centroid estimates. Taking
\begin{equation}\label{eq:N}
 N=\left\lceil8J_{\max}T^2(T^2+4e^2d)\right\rceil
\end{equation}
and applying a conditional union bound gives failure probability at most
$1/(8T^2)$. This reasoning remains valid when each density is selected
adaptively from previous oracle responses and gradients: the moment
bound holds conditional on every prior state, using fresh samples.

For approximate sampling, require each draw, conditional on its current
state and previous draws, to be within total variation
\[
 \xi\le\frac{1}{8T^2J_{\max}N}
\]
of the intended exact law. Sequential maximal coupling makes all samples
coincide with corresponding independent exact draws except on an event
of probability at most $J_{\max}N\xi$. Together with the preceding
centroid bound, this is at most $1/(4T^2)$.
This coupling argument is needed because the Mahalanobis statistic is
not uniformly bounded across states; total variation alone cannot be
applied directly to its second moment.

\begin{thebibliography}{9}

\bibitem{aggarwal2026}
Vaneet Aggarwal.
\newblock Sharp oracle-regret tradeoffs for projection-free online convex optimization.
\newblock arXiv preprint arXiv:2610.00254, 2026.
\newblock \url{https://arxiv.org/abs/2610.00254}.

\bibitem{garberkretzu2022}
Dan Garber and Ben Kretzu.
\newblock New projection-free algorithms for online convex optimization with
adaptive regret guarantees.
\newblock In \emph{Proceedings of the 35th Conference on Learning Theory},
volume 178 of \emph{Proceedings of Machine Learning Research},
pages 2326--2359, 2022.
\newblock \url{https://proceedings.mlr.press/v178/garber22a.html}.

\bibitem{lovaszvempala2007}
L\'aszl\'o Lov\'asz and Santosh Vempala.
\newblock The geometry of logconcave functions and sampling algorithms.
\newblock \emph{Random Structures \& Algorithms}, 30(3):307--358, 2007.
\newblock \href{https://doi.org/10.1002/rsa.20135}{doi:10.1002/rsa.20135}.

\bibitem{mhammedi2022membership}
Zakaria Mhammedi.
\newblock Efficient projection-free online convex optimization with membership oracle.
\newblock In \emph{Proceedings of the 35th Conference on Learning Theory},
volume 178 of \emph{Proceedings of Machine Learning Research},
pages 5314--5390, 2022.
\newblock \url{https://proceedings.mlr.press/v178/mhammedi22a.html}.

\bibitem{mhammedi2025sep}
Zakaria Mhammedi.
\newblock Online convex optimization with a separation oracle.
\newblock In \emph{Proceedings of the 38th Conference on Learning Theory},
volume 291 of \emph{Proceedings of Machine Learning Research},
pages 4033--4077, 2025.
\newblock \url{https://proceedings.mlr.press/v291/mhammedi25a.html}.

\bibitem{mhammedi2025ellipsoid}
Zakaria Mhammedi.
\newblock An ellipsoid algorithm for online convex optimization.
\newblock In \emph{Advances in Neural Information Processing Systems},
volume 38, 2025.
\newblock \href{https://doi.org/10.52202/085713-0030}{doi:10.52202/085713-0030}.

\bibitem{zinkevich2003}
Martin Zinkevich.
\newblock Online convex programming and generalized infinitesimal gradient ascent.
\newblock In \emph{Proceedings of the 20th International Conference on Machine
Learning}, pages 928--936, 2003.
\newblock \url{https://martin.zinkevich.org/publications/ICML03.pdf}.

\end{thebibliography}
\end{document}
